\abstract{
Spatial genetic structure -- the observation that individuals that live closer together are
more genetically similar, termed "isolation by distance" -- is ubiquitious in natural populations.
Despite its prevalence and importance for understanding how genetic variation 
spreads and is shared through populations, much of our contemporary understanding of evolutionary processes
is still informed by models that do not account for this spatial population structure.
While this structure can obscure signatures of evolutionary processes as we understand them,
it also provides another dimension of information that can inform how we analyze genetic variation and
uncover the evolutionary histories that underlie it.

In this dissertation, I use a combination of continuous-space simulation and spatially-referenced
genomic datasets in order to investigate and uncover the evolutionary dynamics of 
\textit{Anopheles} malaria vectors and their \textit{Plasmodium} parasites. 
With this work, I develop and introduce novel spatial methods for detecting evolutionary processes,
including migration and positive selection.
In application, these approaches uncover previously unobserved co-geographic structure
shared by the malaria parasite and its vector, identify novel candidate alleles undergoing positive selection,
and reveal convergent selection pressures acting on deeply-diverged vector species.

In Chapter II, I use machine learning to evaluate evidence for co-geography in the \textit{Anopheles-Plasmodium}
host-parasite system. Interrogating the spatial patterns of prediction errors from a machine learning method 
for geographic prediction from genetic data, I develop a method to reduce prediction bias when training on 
spatially-imbalanced datasets and then use this approach to confirm that errors from both \textit{Anopheles gambiae}
and \textit{Plasmodium falciparum} are a biological signal of shared dispersal patterns.
In Chapter III, I use continuous-space simulation to uncover a novel spatial signal of selective sweeps, and 
use this observation to develop an approach for identifying candidate sweeping alleles in spatially-referenced genetic data.
Applying this method to \textit{Anopheles gambiae} populations in West Africa, I identify several novel alleles
that appear to be under positive selection, including variants with functional relevance and at loci associated with insecticide resistance.
Finally, in Chapter IV, I use a comparative population genetic approach to investigate the evolutionary forces 
shaping and mainintaining patterns of high polymorphism in \textit{Anopheles} vectors. This work reveals
common selection pressures -- primarily acting on genes relevant to immunity and insecticide resistance -- 
that both maintain functional variation within and drive adaptive divergence between species;
in uncovering these patterns, I identify spatially convergent selection on proteins involved in insecticide resistance
across two major African \textit{Anopheles} vectors.
Altogether, this dissertation demonstrates the potential of leveraging spatial population genetic structure in order to 
inform our understanding of evolutionary dynamics, providing insight into the evolutionary history of \textit{Anopheles}
malaria vectors and their \textit{Plasmodium} parasites while introducing novel methods to identify these processes.

This disseration includes previously published and unpublished co-authored material.
}
